\chapter{Horizons, Equation of State, and Vacuum Energy}

\section{Preferred Frame, Superluminal Recession, and the Horizon}

The cosmic microwave background does pick out a special frame. In the frame where that radiation looks isotropic, the nearby galaxies are also, on average, at rest. If we move through the background, photons in the forward direction are Doppler shifted to higher energy and photons in the backward direction to lower energy. The preferred frame is the one in which the microwave background looks isotropic, and it is also, on average, the frame in which the local galaxies are at rest.

\begin{classroomqa}
\lecturetimestamp{00:00:25}
\audiencequestion{Does the apparent wavelength of the cosmic radiation depend on the observer's motion, and does that imply a preferred frame?}
\lecturerresponse{Yes. Motion through the microwave background produces a Doppler anisotropy, so there is a preferred cosmological frame: the one in which the microwave background looks isotropic. It is also, on average, the frame in which the local galaxies are at rest.}
\end{classroomqa}

A second question concerns recession speeds. Hubble's law says that sufficiently distant galaxies separate according to
\begin{equation}
v = H d .
\end{equation}
For large enough \(d\), this recession speed exceeds the speed of light. That does not contradict relativity, because the prohibition is local: one object cannot pass another right in front of us faster than light. The statement about distant galaxies is a statement about the geometry between widely separated places.

\begin{classroomqa}
\lecturetimestamp{00:04:48}
\audiencequestion{Can two galaxies have a relative velocity greater than \(c\)?}
\lecturerresponse{Yes, if they are far enough apart. Hubble recession can exceed the speed of light without violating local relativity, and the place where the recession speed reaches \(c\) is a horizon.}
\end{classroomqa}

Once recession reaches the speed of light, signals sent toward us no longer make net progress. The ray still moves locally at the limiting speed, but the distance it must overcome grows too quickly.

\begin{classroomqa}
\lecturetimestamp{00:07:03}
\audiencequestion{Wouldn't a light ray sent in the reverse direction still travel toward us?}
\lecturerresponse{Locally it does, but space grows faster than the ray can overcome. The ray moves at the limiting speed relative to nearby points, while the total distance back to us keeps increasing.}
\end{classroomqa}

\section{The Expanding Metric and the Line Analogy}

A useful picture is one-dimensional. Imagine a line made of atoms. If we stretch an ordinary material line, eventually it breaks. Now modify the picture: every time neighboring atoms separate too much, we insert another one between them. The line keeps growing, but its local character does not thin out.

\begin{figure}
\centering
\includegraphics[width=0.60\textwidth]{lecture_04_figure_02.png}
\caption{One-dimensional line used to picture an expanding metric.}
\end{figure}

In this picture there is still a local limiting signal speed, just as sound on an ordinary line has a limiting speed. The rule is local: a disturbance cannot outrun that speed relative to nearby points. Nevertheless, if the line is long enough and new points are inserted rapidly enough, sufficiently distant parts of it recede so fast that a signal can never get back.

Special-relativistic velocity addition is therefore only part of the story. It applies locally, when observers pass one another at the same point. It is not the right tool for widely separated galaxies in an expanding geometry.

\begin{classroomqa}
\lecturetimestamp{00:09:42}
\audiencequestion{Does the ordinary relativistic law for adding velocities still hold in this situation?}
\lecturerresponse{Locally, at a point, yes. If observers pass one another one after another, the usual special-relativistic rule applies. But it is not the right way to describe widely separated recession faster than \(c\); there one has to go back to the geometry of expanding space.}
\end{classroomqa}

The same picture gives a simple account of redshift. Draw a wave on the line. As the line expands, the wavelength expands with it. In an expanding geometry this is not merely a suggestive image. Wave equations in expanding spacetime really do produce stretched wavelengths.

\begin{classroomqa}
\lecturetimestamp{00:14:42}
\audiencequestion{Is the statement that a wave stretches with the universe theoretical, empirical, or both?}
\lecturerresponse{It is a theoretical consequence of wave propagation in expanding spacetime, and it is also supported by astronomical observation.}
\end{classroomqa}

This picture is not a rubber band. A rubber band thins out when stretched, and someone living in such a medium would see the local properties change. Here the local properties stay the same while the global distance scale grows.

\section{Last Scattering and the Observable Universe}

The horizon is the place where recession reaches the speed of light. The surface of last scattering lies slightly inside it. A signal emitted from the horizon would be stretched without bound; the last-scattering radiation has a huge but finite redshift.

\begin{classroomqa}
\lecturetimestamp{00:17:53}
\audiencequestion{What is the relationship between superluminal recession, the surface of last scattering, and the observable universe?}
\lecturerresponse{The horizon is where recession reaches the speed of light. The surface of last scattering lies slightly inside it, because its redshift is enormous but finite. For practical purposes it marks the edge of the observable universe.}
\end{classroomqa}

The redshift of the last-scattering radiation is roughly a factor of \(10^3\). That is already so large that the surface of last scattering is, for many purposes, very close to the horizon. The observable universe is the portion of the universe from which light can reach us.

At late times the expansion approaches the vacuum-dominated regime. Then the Hubble quantity tends toward a constant, the horizon approaches a fixed size, and both galaxies and microwave radiation are progressively lost from view.

\begin{classroomqa}
\lecturetimestamp{00:22:43}
\audiencequestion{Since \(H\) changes with time, is the horizon distance changing, and will the microwave background eventually disappear?}
\lecturerresponse{At late times \(H\) tends to a constant, so the horizon approaches a fixed distance. The microwave background keeps stretching, cooling, and fading away. Distant galaxies also pass through the horizon and disappear from view.}
\end{classroomqa}

Only structures that are already gravitationally bound remain visible. The solar system does not fly apart, and the galaxy does not fly apart, but the large-scale universe becomes an isolated island.

\section{The Basic Expansion Equation and the Three Components}

Now the basic expansion equation is
\begin{equation}
\left(\frac{\dot a}{a}\right)^2 = \frac{8\pi G}{3}\rho ,
\qquad c=1 .
\end{equation}
Here \(a(t)\) is the scale factor, and \(\rho\) is the total energy density.

\begin{figure}
\centering
\includegraphics[width=0.74\textwidth]{lecture_04_figure_03.png}
\caption{Basic expansion equation together with the one-dimensional line analogy.}
\end{figure}

There are three dynamically important components.

For matter,
\begin{equation}
\rho=\frac{\text{const.}}{a^3},
\qquad
a(t)\sim t^{2/3}.
\end{equation}
This is the case of massive objects moving with the overall expansion and essentially at rest in the comoving frame.

For radiation,
\begin{equation}
\rho=\frac{\text{const.}}{a^4},
\qquad
a(t)\sim t^{1/2}.
\end{equation}
The extra power of \(a\) comes from the redshifting of each photon as the universe expands.

There is also
\begin{equation}
\rho_0=\text{constant},
\end{equation}
the vacuum energy, also called dark energy or, up to conventional numerical factors, the cosmological constant. This component does not dilute at all.

The next step is to relate energy density to pressure.

\section{Radiation Pressure and Cold Matter}

Start with radiation. Consider a one-dimensional box, just a line segment of length \(L\), containing photons bouncing back and forth. In one dimension, force and pressure are effectively the same thing.

Take a photon of momentum magnitude \(p\). When it reflects from a wall, its momentum changes by
\begin{equation}
\Delta p = 2p.
\end{equation}
The time between successive collisions with the same wall is
\begin{equation}
\Delta t = \frac{2L}{c}.
\end{equation}
Therefore the average force on the wall is
\begin{equation}
F=\frac{\Delta p}{\Delta t}=\frac{pc}{L}.
\end{equation}
For a photon,
\begin{equation}
E=pc,
\end{equation}
so
\begin{equation}
F=\frac{E}{L}.
\end{equation}
The energy divided by the length is the one-dimensional energy density. In this one-dimensional case,
\begin{equation}
P=\rho .
\end{equation}

With many photons the same relation survives after averaging. The total force is just the total energy per unit length.

Now go to three dimensions. If all photons moved along one axis only, they would exert pressure on one pair of faces and none on the others. That is deliberately wrong, but it shows immediately what has to be corrected. In an isotropic bath, momentum is distributed over all directions. On average, one third of the energy density contributes to each pressure component. The result is
\begin{equation}
P=\frac{\rho}{3}.
\end{equation}

For cold matter the answer is different. Take a comoving box, a small region moving with the flow of the universe. The particles inside carry energy, but if they are at rest relative to the box they do not bang into its walls. The pressure is therefore
\begin{equation}
P=0.
\end{equation}

\begin{figure}
\centering
\includegraphics[width=0.74\textwidth]{lecture_04_figure_04.png}
\caption{Matter and radiation: \(a\sim t^{2/3}\), \(P=0\), and \(a\sim t^{1/2}\), \(P=\rho/3\).}
\end{figure}

The board summary is
\begin{align}
\text{matter:}\qquad & a(t)\sim t^{2/3}, \qquad P=0,\\
\text{radiation:}\qquad & a(t)\sim t^{1/2}, \qquad P=\frac{\rho}{3}.
\end{align}

\begin{classroomqa}
\lecturetimestamp{00:54:27}
\audiencequestion{In the expanding universe, aren't the particles actually moving apart?}
\lecturerresponse{Yes, but the pressure argument is applied to an infinitesimal comoving box. On that local scale the thermodynamic reasoning still works.}
\end{classroomqa}

\section{The Equation of State}

Pressure is taken to be proportional to energy density,
\begin{equation}
P=w\rho .
\end{equation}
On the board the same parameter is written as \(\omega\). For the two basic cases,
\begin{equation}
w=0 \quad \text{for cold matter},
\qquad
w=\frac{1}{3} \quad \text{for radiation}.
\end{equation}

Intermediate cases are possible. Values between \(0\) and \(1/3\) describe material that is neither completely cold nor fully relativistic. Warm dark matter is the example mentioned here.

\begin{figure}
\centering
\includegraphics[width=0.74\textwidth]{lecture_04_figure_05.png}
\caption{Equation of state \(P=\omega\rho\) together with the density-scaling argument.}
\end{figure}

The next question is whether \(w\) controls the dependence of \(\rho\) on \(a\). The answer comes from a thermodynamic identity.

Take a box of volume \(V\). If the contents exert pressure \(P\), then when the box expands by \(dV\) the work done on the walls is
\begin{equation}
dE=-P\,dV.
\end{equation}
Now write the energy as
\begin{equation}
E=\rho V.
\end{equation}
Its differential is
\begin{equation}
dE=\rho\,dV+V\,d\rho .
\end{equation}
Equating the two expressions gives
\begin{equation}
\rho\,dV+V\,d\rho=-P\,dV,
\end{equation}
or
\begin{equation}
V\,d\rho=-(\rho+P)\,dV.
\end{equation}
Using \(P=w\rho\),
\begin{equation}
V\,d\rho=-(1+w)\rho\,dV.
\end{equation}
Divide by \(\rho V\):
\begin{equation}
\frac{d\rho}{\rho}=-(1+w)\frac{dV}{V}.
\end{equation}
Integrating,
\begin{equation}
\rho \propto V^{-(1+w)}.
\end{equation}
For a comoving cosmological box,
\begin{equation}
V\propto a^3,
\end{equation}
and therefore
\begin{equation}
\rho \propto a^{-3(1+w)}.
\end{equation}

The earlier cases are recovered at once:
\begin{align}
w=0 &\Longrightarrow \rho \propto a^{-3},\\
w=\frac{1}{3} &\Longrightarrow \rho \propto a^{-4}.
\end{align}

This is why \(w\) matters. Once \(\rho(a)\) is known, it can be put back into the Friedmann equation. One then tries a power law \(a\sim t^\alpha\) and determines how the universe expands for that equation of state.

\section{Vacuum Energy and Negative Pressure}

One special case behaves differently:
\begin{equation}
w=-1.
\end{equation}
Then
\begin{equation}
1+w=0,
\end{equation}
so the exponent vanishes and
\begin{equation}
\rho=\rho_0=\text{constant}.
\end{equation}
The corresponding pressure is
\begin{equation}
P=-\rho .
\end{equation}

Negative pressure is not mysterious in itself. Positive pressure pushes outward on the walls of a box. Negative pressure pulls inward. It is tension. A stretched rubber band gives the everyday example: it stores positive energy, but it pulls its ends together.

Substitute the constant density \(\rho_0\) into the Friedmann equation:
\begin{equation}
\left(\frac{\dot a}{a}\right)^2=\frac{8\pi G}{3}\rho_0.
\end{equation}
The right-hand side is just a number, so
\begin{equation}
\frac{\dot a}{a}=\sqrt{\frac{8\pi G}{3}\rho_0}.
\end{equation}
This is the special case in which the Hubble quantity really is constant. The equation can be written as
\begin{equation}
\dot a=\sqrt{\frac{8\pi G}{3}\rho_0}\,a,
\end{equation}
and the solution is
\begin{equation}
a(t)\propto \exp\!\left(\sqrt{\frac{8\pi G}{3}\rho_0}\,t\right).
\end{equation}

This is exponential expansion. The relevant timescale is the \(e\)-folding time, not a doubling time, and the rough estimate used here puts it on the order of the age of the universe. Vacuum energy and cosmological constant are the same physical ingredient apart from historical numerical conventions.

\begin{classroomqa}
\lecturetimestamp{01:27:16}
\audiencequestion{If vacuum energy drives exponential expansion, won't everything eventually be ripped apart?}
\lecturerresponse{Not necessarily. Bound systems can resist it. A galaxy, a solar system, or an atom need not be torn apart, whereas sufficiently large and dilute structures do participate in the expansion and eventually separate.}
\end{classroomqa}

\begin{classroomqa}
\lecturetimestamp{01:31:43}
\audiencequestion{Is there anything interesting in the negative sign when one takes the square root of the Friedmann equation?}
\lecturerresponse{Yes. The square root gives two branches. One is the expanding solution, and the other is the time-reversed contracting solution. The equations themselves admit both.}
\end{classroomqa}

\section{Three Eras and the Observed Expansion History}

All three components are present, and they scale differently:
\begin{equation}
\begin{aligned}
\rho_{\text{radiation}}&\propto a^{-4},\\
\rho_{\text{matter}}&\propto a^{-3},\\
\rho_0&=\text{constant}.
\end{aligned}
\end{equation}
That immediately implies three eras.

At very small \(a\), radiation dominates. Later, matter dominates. At sufficiently large \(a\), vacuum energy dominates. The corresponding expansion laws are
\begin{equation}
a(t)\sim t^{1/2},
\qquad
a(t)\sim t^{2/3},
\qquad
a(t)\sim e^{Ht}.
\end{equation}

The rough times quoted here are these: radiation gives way to matter after about \(10^4\) years. The universe becomes transparent later, after about \(10^5\) years. The transition from matter domination to accelerated expansion occurs much later, after several billion years, roughly halfway through cosmic history in the order-of-magnitude estimate being used.

How is this history inferred? We cannot see the future, but we can see the past. Looking to greater distance means looking to earlier time. By observing distant galaxies, measuring how their motion correlates with distance, and reconstructing the Hubble quantity as a function of lookback time, one reconstructs the time history of the scale factor. The radiation-dominated era is too early to be seen directly in ordinary galaxy astronomy, but the matter-dominated era and the transition to vacuum domination are visible in the reconstructed curve.

\section{Last Scattering in the Far Future}

The future microwave sky is not the same story as galaxies disappearing through the horizon. The radiation we see at later times comes from farther matter, so the effect is seeing deeper and cooler rather than watching the same atoms simply recede away.

\begin{classroomqa}
\lecturetimestamp{01:48:07}
\audiencequestion{If the surface of last scattering fades into the horizon, is that the same thing as galaxies moving beyond the horizon?}
\lecturerresponse{No. The later radiation does not come from the same atoms as the radiation we see now. It comes from atoms that were farther away, so it is like seeing deeper and deeper into the universe, while the radiation looks cooler and cooler.}
\end{classroomqa}

\section{Curvature, the Big Crunch Picture, and the Late-Time Horizon}

Up to this point the curvature term has been suppressed. Restoring it gives
\begin{equation}
\left(\frac{\dot a}{a}\right)^2
=
\frac{8\pi G}{3}\rho
-
\frac{K}{a^2}.
\end{equation}
Numerically this term seems small, but conceptually it matters. In the older Newtonian language it distinguishes the cases of positive, zero, and negative total energy. In the relativistic language it distinguishes the spatial geometries.

Before vacuum energy was taken seriously, one often imagined a matter-dominated universe in which the expansion would eventually reverse. The result was the old Big-Crunch picture: expansion from the Big Bang, then recontraction. The mathematics is still valid, but it no longer appears to describe our universe.

\begin{classroomqa}
\lecturetimestamp{01:53:51}
\audiencequestion{Is this vacuum energy related to quantum vacuum energy?}
\lecturerresponse{Yes, or at least it is believed to be.}
\end{classroomqa}

The rough composition today is
\begin{equation}
\frac{\rho_0}{\rho_{\text{total}}}\sim 0.7,
\qquad
\frac{\rho_{\text{matter}}}{\rho_{\text{total}}}\sim 0.3.
\end{equation}
Of that matter component, luminous matter is about \(0.05\) and dark matter about \(0.25\) of the total. Radiation is negligible today. Dark and luminous matter scale in the same way, so their ratio stays fixed even while the relative abundance of matter, radiation, and vacuum changes from era to era.

In the vacuum-dominated limit,
\begin{equation}
H=\sqrt{\frac{8\pi G}{3}\rho_0}.
\end{equation}
Then Hubble's law becomes
\begin{equation}
v=Hd.
\end{equation}
The horizon is defined by \(v=1\), so its distance is
\begin{equation}
d_H=\frac{1}{H}
=
\frac{1}{\sqrt{8\pi G\rho_0/3}}.
\end{equation}
The rough number given here is about \(10\) to \(12\) billion light years.

The important point is the late-time behavior. Once the expansion is exponential, the horizon distance becomes fixed. Galaxies continue to pass through it and disappear by becoming more and more redshifted. The microwave background is redshifted away as well. In the remote future, only the nearby gravitationally bound structures remain visible, and the universe looks like an isolated island.

\section{Geometry from the Last-Scattering Surface}

Finally the geometry can be tested from the last-scattering surface. Suppose the distance to that surface is known. Suppose also that one can identify a lump there whose physical size is known well enough. Then one can compare that size and distance with the observed angular size.

In flat space the angle is the Euclidean one. In positively curved space it would appear larger; in negatively curved space, smaller. The analogy is ordinary surveying: take an object of known size at known distance and compare the measured angle with Euclidean geometry.

\begin{classroomqa}
\lecturetimestamp{02:08:09}
\audiencequestion{What is an example of an object out there whose size we know?}
\lecturerresponse{The largest lumps in the microwave background. We think we know their size from how long they took to evolve and from the speed of sound.}
\end{classroomqa}

The relevant objects here are the largest lumps on the microwave background. Their observed angular size comes out consistent with Euclidean geometry to within a few percent. On that basis, space appears flat.

\section{Measured \(w\) and a Final Question}

The observational importance of \(w\) is now clear. If astronomers report that \(w\) is very close to \(-1\), they are saying that the vacuum-energy component is very nearly constant. That is exactly the point made here.

That does not mean general relativity forces \(w\) to be exactly \(-1\). The point here is empirical. The lecture briefly notes that more radical possibilities are not excluded in principle, but that the nearly constant vacuum-energy case is the one for which there is good reason to believe.

\begin{classroomqa}
\lecturetimestamp{02:11:04}
\audiencequestion{In the density calculation for matter and for radiation, did we assume that the energy was constant and then adjust the scale factor to figure out the formula? And if that is true for the first two cases, what happens in the vacuum-energy case, where the density stays constant?}
\lecturerresponse{No. I did not assume it was constant. We used that the change in energy is \(-P\,dV\). For an imaginary comoving box moving along with the expansion, work is done on the walls, and that accounts for the change in energy. The vacuum-energy case raises the question in a different way, and the closing remark here is that energy-conservation language in cosmology does not say very much, because the total energy of the universe is zero.}
\end{classroomqa}
